\clearpage
\section{Conclusions}
\label{sec:conclusions-final}
\begingroup
\fontsize{10.2}{12.2}\selectfont
\setlength{\parskip}{2.5pt plus .5pt}

La thèse géométrique est établie par reconstruction et transport : une représentation positive de la structure discrète du continu conserve, avec ses marques, le registre dont elle provient et les lecteurs que celui-ci détermine. La charge totale et les quotients de mineurs restituent chaque coordonnée ; les deux mémoires nonadiques fournissent une autre reconstruction opératorielle. La valeur et la forme s'articulent ainsi par des inverses explicites, après la génération APP--TRIT--TPK et ses prolongements compatibles. La double projection maintient unies la réalisation arithmétique et l'incidence exceptionnelle dans cette généalogie.

Le raffinement dépasse la limitation d'une matrice à treize colonnes. Pour tout rang et toute dimension externe finis, un niveau nonadique suffisant produit les mineurs positifs et l'image amplituèdrale correspondante. Les nœuds hérités conservent leurs drapeaux ; les moyennes cellulaires conservent la mesure et enregistrent le détail orthogonal. En rang un, la couverture simpliciale et la récurrence des résidus s'étendent à toute dimension finie. Dans les rangs supérieurs, le résultat détermine des collections positives compatibles, de rang et de limite démontrés ; leur étendue dimensionnelle demeure distincte de la classification globale de toutes les cellules amplituèdrales.

La conservation admet trois expressions complémentaires. La forme canonique s'additionne sous subdivision d'un support fixe ; l'énergie effective est conservée par réduction de Schur ; et la mesure chronologique détermine des moments dont l'incrément se décompose en détails récupérables. Ces moments construisent des opérateurs de Jacobi dont les compressions finies ont un spectre simple et des résidus positifs. On obtient ainsi un parcours effectif des séparations arithmétiques à une représentation spectrale, avec une fraction continue déterminée par la même mesure. Les mutations qui ne changent que la coordinatisation conservent les lecteurs du point marqué.

La réalisation physique gagne en précision grâce à ses opérateurs. Les hamiltoniens des trois régimes tritiques possèdent des formes symplectiques et des transformations de Legendre explicites. La coordinatisation elliptique variable incorpore son terme de connexion ; la mémoire parabolique conserve son invariant ; l'ouverture hyperbolique possède son évolution propre. Les lecteurs d'action et de réponse constitutive sont transportés par l'inverse grassmannien. Avec les représentations régionales et les marques entières déclarées, la réponse constitutive étiquetée restitue à son tour le registre. La condition d'égalité des traces gouverne le recollement des formes à valeurs dans la droite d'action.

La déformation exceptionnelle conserve le volume et la symétrie
ternaire, et transforme l’inversion du paramètre en dualité
réticulaire. Sa réciprocité thêta et sa réalisation de signature
déployée relient la variation géométrique à la dualité du
réseau. En Yang--Mills, l’échelle aréale--volumétrique, la mesure commune,
la coercivité uniforme et le témoin persistant situent l’écart
auprès du Hamiltonien qui le réalise. Le développement conserve
l’espace physique et les opérateurs d’entrelacement nécessaires pour transporter
cette assertion à travers le raffinement.

Le résultat conjoint est une architecture de formes avec mémoire : la dimension peut augmenter, la représentation peut changer et l'information pertinente demeure récupérable au moyen d'applications déterminées. Tel est le contenu distinctif de l'article. La géométrie positive, la dualité réticulaire et la réalisation spectrale sont reliées par des constructions qui spécifient ce qu'elles engendrent, ce qu'elles conservent et comment elles se reconstruisent à partir d'une même origine discrète.
La représentation d'incidence conserve également la composante uniforme du registre : son inverse reconstruit les douze coordonnées et permet de récupérer la direction géométrique après le transport de mémoire. Dans la réalisation sectorielle, l'aire totale et l'occupation pondérée déterminent conjointement la répartition hexade--octade. La purification spécifiée transforme cette répartition en une structure d'intrication, avec une identité finie entre entropie, aires sectorielles et entropie relative. La conservation de l'information est ainsi précisée pour chaque application : registre complet, direction, aire et composition possèdent des inverses ou des fibres explicitement distincts.
L'échelle d'aire est déterminée au moyen d'une longueur engendrée par la composition linéaire du retour enregistré et de la norme locale de Hadamard. La réciprocité radiale détermine le coefficient gravitationnel dans la réalisation réduite déclarée et reconnaît ensuite cette longueur comme la longueur de Planck. Les mutations du même point marqué et les raffinements hérités conservent l'opérateur d'aire, ses multiplicités et l'identité modulaire. Les formes canoniques à valeurs dans cet opérateur étendent la conservation aux résidus, tandis que le rayon de chaque horizon maintient son propre facteur géométrique. Cette composition explicite la dépendance entre longueur, incidence et aire sans altérer les états sectoriels.
\par\endgroup
\clearpage
